\section{Verifier-supervised look-back}
\label{sec:method}
Both instantiations learn which earlier screen to re-read from the same verifier output. They differ in how the history reaches the model, which also means the same 188 cases become 2,861 training items for one and 188 for the other.

\paragraph{Labels and pairs.} Each verified case $(e,t)$ yields one training item per earlier screen $i<t$. For the per-screen scorer the item is a pair. The input holds the task description, the OCR text of screen $t$, the OCR text of screen $i$, and the distance $t-i$; the label is 1 if $i$ is in the verifier's set and 0 otherwise. The 188 cases yield 2,861 pairs, 463 of them positive. For the full-history reader the item is one example per case: the input holds the task and the OCR text of every earlier screen, and the target is the number of the latest screen in the verifier's set. No human labels any item.

\paragraph{Per-screen scorer.} Qwen2.5-7B-Instruct~\citep{qwen25} with a LoRA (low-rank adaptation) adapter of rank 16 on the attention projections (40M trainable parameters). The prompt ends with ``Answer Yes or No''; we take the logits of \texttt{Yes} and \texttt{No} at the answer position, apply a two-way softmax, and minimize cross-entropy with positives weighted 3 to 1. Localization is the argmax of $P(\text{Yes})$ over earlier screens, computed by $t$ forward passes of 600--900 tokens each and no generation.

\paragraph{Full-history reader.} The same base model and adapter rank, fine-tuned to emit the screen number from a prompt that holds the whole history, with the loss on the number only. Localization is one greedy generation of at most five tokens. Both variants train for 2 epochs at learning rate $10^{-4}$ with 4 items per update, 13--27 minutes per fold on one 48GB GPU.

\paragraph{Cost.} At $t=16$ the scorer reads about 12K tokens across $t$ small forward passes; the full-history reader reads 1--5K tokens in one pass, and prompted full history reads the same. A 3B scorer loses 4.2 points overall and gains 11.1 at $t\ge20$ (\S\ref{sec:ablation}), so the scorer's cost is adjustable.

\paragraph{Read stage.} After localization, a Qwen2.5-14B-Instruct reader types the text, with the prompt shared by every condition. \emph{Single-screen read} receives the task, the current screen, and the chosen screen's OCR text. \emph{Read with history} receives the full history as well, with the chosen screen repeated and marked before the current screen; it costs full history's context plus one repeated screen.

\paragraph{Protocol.} We split the 188 cases into five folds, keeping every case from one trajectory in the same fold so that no trajectory appears in both train and test. We train on four folds, score the fifth, and pool the five scored folds; every cross-validated number below is that pool. For the external set we train once on all 188 cases and evaluate once. \S\ref{sec:index} states which writer numbers are not cross-validated.
